% TOP_PROBLEM: {"id":30007654,"problem_number":"LOCAL-30007654","title":"Cash transfers at economy-wide scale","category_id":21,"category":{"id":21,"name":"economics","display_name":"Economics","description":"Open economic theory statements, published research agendas, and proposed scoped specifications; question provenance and historical priority are qualified per record.","slug":"economics","order_index":21,"created_at":"2026-10-08T00:00:00Z"},"status":"provenance_pending","external_url":null,"legacy_ids":[],"published":false,"statement":"In the explicitly specified setting: When do large cash-transfer programs create persistent productivity gains and beneficial spillovers, and when do inflation, financing costs, or supply constraints erode them? The mathematical statement above fixes the target, assumptions and scope of this catalogue formulation.","economics":{"original_id":143,"field":"Development and poverty","kind":"empirical","formalization_basis":"proposed_specification","source_status":"not_verified","priority_status":"not_established","priority_note":"No qualifying problem statement was directly verified, so historical priority cannot be assigned.","earliest_verified_reference":null,"current_reference":{"authors":["Tommaso Crosta","Dean Karlan","Finley Ong","Julius Rüschenpöhler","Christopher Udry"],"title":"Unconditional Cash Transfers: A Bayesian Meta-Analysis of Randomized Evaluations in Low and Middle Income Countries","year":2024,"url":"https://www.nber.org/papers/w32779","doi":"10.3386/w32779","locator":"Abstract, August 2026 revision, paragraphs on dynamic effects and context heterogeneity","evidence_summary":"The meta-analysis provides positive effects and investigates heterogeneity; it does not explicitly state the proposed equilibrium-and-scale question as open."},"formalization_note":"This is a proposed scoped research specification. No primary passage explicitly stating this entire remaining question as open was verified. The supporting literature may answer important subquestions; this entry must not be advertised as a verified formal open problem."},"statusQualification":"Provenance pending: an explicit published statement of this exact open question was not verified. This is a catalogue-authored candidate specification, not a certified open conjecture. This is a proposed scoped research specification. No primary passage explicitly stating this entire remaining question as open was verified. The supporting literature may answer important subquestions; this entry must not be advertised as a verified formal open problem.","statusReviewedAt":"2026-10-08"}
% FIRST_POSTED: 2026-10-08
% ENTRY_KIND: statement_only
% !TeX program = lualatex
\documentclass[11pt]{article}
\usepackage{fontspec}
\usepackage[a4paper,margin=25mm]{geometry}
\usepackage{amsmath,amssymb,mathtools}
\usepackage{xurl}
\usepackage[hidelinks]{hyperref}
\newcommand{\sref}[2]{\hyperlink{source:#1:#2}{[S#2]}}
\setlength{\emergencystretch}{3em}
\begin{document}
% problemId:30007654
\section*{Cash transfers at economy-wide scale}\label{30007654}
\subsection{Definitions and mathematical statement}
Consider a specified national or regional economy with eligible and ineligible households, firms, and government. Policy \(p=(b,q,f)\) gives transfer amount \(b\) to coverage share \(q\) and uses financing rule \(f\), which explicitly specifies taxes, borrowing, aid, or spending displacement. Let \(C_{it}(p)\) be household real consumption, \(P_t(p)\) a price index, and \(Y_t(p)\) real aggregate production. Define \(\tau_C(T,p)=E[\sum_i\omega_i\{C_{iT}(p)-C_{iT}(0)\}]\) for declared population weights \(\omega_i\), with analogous price and output responses. Estimate how these effects vary across transfer intensity, supply capacity, financing, and time since disbursement. Include effects on nonrecipients, migration, firm entry, and government budgets. Randomized household transfers identify local effects only under an appropriate interference design; extrapolation to full coverage requires a validated equilibrium model or policy variation at larger scales. Determine which productivity and consumption changes survive after transfers cease. The target is a family of context-specific response functions, not one universally positive cash-transfer multiplier. Existing evidence of particular general-equilibrium effects should be treated as partial answers; an explicit source declaring this whole catalogue formulation open was not verified.

\par\textbf{Formulation and scope.} This is a proposed scoped research specification. No primary passage explicitly stating this entire remaining question as open was verified. The supporting literature may answer important subquestions; this entry must not be advertised as a verified formal open problem.

\par\textbf{Open-question provenance.} An explicit published open-question statement was not verified. Sources below are supporting literature and must not be treated as a first listing.

\par\textbf{Historical priority.} No qualifying problem statement was directly verified, so historical priority cannot be assigned.

\subsection{Short English statement}
In the explicitly specified setting: When do large cash-transfer programs create persistent productivity gains and beneficial spillovers, and when do inflation, financing costs, or supply constraints erode them? The mathematical statement above fixes the target, assumptions and scope of this catalogue formulation.

\subsection{Sources}
\begin{itemize}\raggedright
\item[\textnormal{[S1]}]\hypertarget{source:30007654:1}{} Tommaso Crosta, Dean Karlan, Finley Ong, Julius Rüschenpöhler, Christopher Udry. 2024. Unconditional Cash Transfers: A Bayesian Meta-Analysis of Randomized Evaluations in Low and Middle Income Countries. Abstract, August 2026 revision, paragraphs on dynamic effects and context heterogeneity.
\url{https://www.nber.org/papers/w32779}.
DOI: 10.3386/w32779.
{\small\textit{Supporting or current literature; see evidence qualification. The meta-analysis provides positive effects and investigates heterogeneity; it does not explicitly state the proposed equilibrium-and-scale question as open.}}
\end{itemize}
\end{document}
